\begin{table}[H]
\centering
\sffamily
\renewcommand{\arraystretch}{1.3}
\begin{threeparttable}
\caption{D\'ecomposition de la croissance annuelle moyenne de la productivit\'e du travail (\%)}
\label{tab:note_decomposition}
\begin{tabular}{l*{4}{c}}
\toprule
& 1961--2019 & 1961--1980 & 1980--2000 & \textbf{2000--2019} \\
\midrule
$\Delta \ln(Y/L)$ & 1.28 & 2.01 & 1.32 & \textbf{0.52} \\
\quad PTF intra-industries & 1.28 & 1.79 & 1.26 & \textbf{0.20} \\
\quad R\'eallocation (Baumol) & -0.44 & -0.13 & -0.27 & \textbf{-0.33} \\
\quad Approfondissement du capital & 0.44 & 0.35 & 0.33 & \textbf{0.65} \\
\bottomrule
\end{tabular}
\begin{tablenotes}\footnotesize
\item \textit{Note}\,: Croissance annuelle moyenne en points de pourcentage. Les trois composantes s'additionnent exactement \`a $\Delta \ln(Y/L)$. La PTF intra-industries et la r\'eallocation sont amplifi\'ees par le facteur $1/(1-\bar\alpha) > 1$, qui refl\`ete l'effet indirect de la PTF sur l'accumulation du capital (voir annexe~A). Source\,: Statistique Canada, tableau 36-10-0217-01.
\end{tablenotes}
\end{threeparttable}
\end{table}